\chapter{Discussion}
\label{ch:discussion}

This chapter interprets the experimental results presented in Chapter~\ref{ch:experiments},
examines the counter-intuitive relationship between formal conformance and strict success
rate, draws practical implications for integrating the HSP-Agile pipeline into sprint
workflows, and situates the approach within the broader landscape of related research.
Limitations and threats to external validity are then assessed.

% ---------------------------------------------------------------------------
\section{Interpretation of Core Findings}
\label{sec:disc:interpretation}
% ---------------------------------------------------------------------------

The central empirical finding of this thesis is that the full HSP-Agile pipeline (mode~M)
achieves a mean strict formal conformance of \textbf{90.4\%} on the hard benchmark of 120
precedence-sensitive Agile-SOFL tasks.  This represents a gain of \textbf{6.2 percentage
points} over the one-shot LLM baseline (B1: 84.2\%) and \textbf{2.4 percentage points}
over the test-feedback loop baseline (B2: 88.0\%).

Two interlocking mechanisms produce this improvement.  First, the formal conformance checker
generates Z3-driven test cases targeted specifically at the ordering constraints encoded in
each Functional Scenario Framework (FSF) specification~\citep{barrett2018smt,demoura2008z3,liu2014sofl}.
These cases expose ordering violations that coarser unit tests, which typically exercise only
the ``happy path'' of each scenario branch, systematically miss~\citep{zhu1997specification,ammann2016testing}.
Second, the counterexample-guided repair loop~\citep{solarlezama2008sketching}
converts witnessed failures into structured repair prompts that direct the language model
towards the violated constraint~\citep{lynn2024verifierloop,chen2021codex}.

The low conformance of the B0 reference mode (27.5\%) merits comment.  B0 consists of
auto-generated reference implementations produced by a template-based generator that
correctly handles straightforward scenario orderings but is not optimised for the
hardest precedence-sensitive patterns in the benchmark.  Its score therefore establishes a
meaningful lower bound: it demonstrates how poorly non-LLM template generation performs
on the hard partition, and it calibrates the difficulty of the benchmark rather than
representing a realistic baseline for production use.

The ablation results confirm that both principal pipeline stages contribute independently.
Removing the formal checker entirely (A1) drops conformance to 86.5\% ($-3.9$\,pp
vs.\ M), while removing the repair loop (A3) drops it to 84.1\% ($-6.3$\,pp vs.\ M).
Removing the pattern guard (A2) yields a marginal gain of $+0.5$\,pp, a finding discussed
in detail in Section~\ref{sec:disc:implications}.  Together, the ablation deltas account
for the bulk of the conformance advantage of M over B1, validating the design rationale of
each pipeline stage.

% ---------------------------------------------------------------------------
\section{The Quality–Pass-Rate Paradox}
\label{sec:disc:paradox}
% ---------------------------------------------------------------------------

One of the most counter-intuitive results in the evaluation is that the full pipeline M
exhibits a \emph{lower} strict success rate (25.0\%) than both B1 and B2 (26.7\%), despite
achieving higher formal conformance.  This apparent paradox requires careful unpacking,
because it has direct consequences for how one should choose evaluation metrics in
formal-conformance studies.

The two metrics measure fundamentally different things.  \emph{Strict success rate} is a
binary per-task metric: a task is marked successful only when \emph{all} generated test
cases pass.  A single failing case, regardless of how close the candidate is to being
correct, yields a score of zero for that task.  \emph{Formal conformance}, by contrast, is
a continuous metric: it records the fraction of test cases that pass, averaged over all
tasks.  A candidate that satisfies 9 out of 10 generated cases contributes 0.9 to the
conformance average even though it contributes 0 to the strict success count.

The paradox arises because M's conjunctive acceptance gate is deliberately conservative.
When a candidate passes all test cases but the pattern guard detects a residual fault
pattern, M classifies the candidate as ``best effort'' and does \emph{not} emit a strict
accept.  In other words, M may \emph{raise the bar for acceptance} relative to B1 and B2,
which accept any candidate that clears their respective stopping criteria.  The result is
that a subset of tasks that B1 or B2 would count as strict successes are counted by M as
non-successes because the pattern guard fired — yet those same candidates contribute fully
to the conformance average, inflating M's conformance relative to its strict success rate.

A useful analogy clarifies this relationship.  Suppose two students sit ten examinations.
Student~A scores 89\% on every paper.  Student~B scores 100\% on three papers and 65\% on
the remaining seven.  On a binary pass/fail threshold of 75\%, Student~A passes all ten and
Student~B passes only three.  Yet Student~B's \emph{mean score} is $0.3 \times 100 + 0.7
\times 65 = 75.5\%$, below Student~A's 89\%.  The binary metric favours Student~B despite
lower average quality; the continuous metric correctly orders them.  The HSP-Agile results
follow the same logic: M is the student who scores consistently high across all tasks
(conformance), whereas B1 and B2 occasionally solve tasks completely but average lower
quality.

This analysis provides a clear empirical justification for adopting formal conformance as
the \emph{primary} metric on the hard benchmark.  The strict success rate is too coarse to
differentiate between a pipeline that produces consistently high-quality candidates and one
that occasionally solves a task while producing low-quality candidates on the majority.
The Wilcoxon signed-rank test on strict success rate yields $p = 0.921$ (not significant),
consistent with the interpretation that the binary metric cannot distinguish M from B2 on
this benchmark; the difference is effectively noise in the binary domain even though a
meaningful quality difference exists in the continuous domain.

% ---------------------------------------------------------------------------
\section{Implications for Agile-Formal Integration}
\label{sec:disc:implications}
% ---------------------------------------------------------------------------

\subsection{Sprint Integration and CI/CD Deployment}

The mean latency of mode~M is 43{,}659\,ms per task, approximately 4.6 times the latency
of B2 (10{,}769\,ms).  This overhead reflects the additional LLM API calls required by the
repair loop (mean 2.08 calls per task under M) and the Z3-based conformance checker.
Whether this latency is acceptable depends critically on the deployment context.

For \emph{interactive} use — a developer waiting in the IDE for a candidate implementation
— 44 seconds is at the upper boundary of acceptable response times and would likely
require UX mitigations such as streaming partial results or asynchronous notification.
However, for \emph{CI/CD pre-merge gates}, 44 seconds per task is well within the range
of typical automated quality checks.  A sprint team merging 20--30 LLM-assisted
implementation units per day would incur roughly 15--22 minutes of additional checking
time in the pipeline, an overhead comparable to compiling a medium-sized Java project or
running a unit-test suite.  Given that the alternative — propagating conformance faults
into downstream Agile-SOFL verification artefacts — can require hours of manual
debugging, the trade-off is clearly favourable in a formal-methods-adjacent workflow.

We therefore recommend deploying HSP-Agile as a background pre-merge quality gate rather
than as an interactive code-completion tool.  This framing positions the 44-second latency
not as a user-experience constraint but as a CI step, a context in which it is unremarkable,
consistent with agile fitness-function practice~\citep{harman2012formalagile,beck2001xp}.

\subsection{The Formal Checker as a Lightweight Oracle}

The Z3-based conformance checker occupies an important middle ground between no formal
verification and full deductive verification.  Full verification tools such as Dafny~\citep{leino2010dafny} and Coq provide complete guarantees but require annotation-heavy specifications, trained
verification engineers, and verification times that can span hours for non-trivial
programs~\citep{woodcock2009formalsurvey}.  The Z3 checker used in HSP-Agile, by contrast, provides \emph{bounded}
guarantees: it guarantees conformance with respect to the finite test suite it generates,
not with respect to all possible inputs.  Within those bounds, the overhead per task is
less than 100\,ms — a negligible fraction of the total pipeline latency.

This lightweight oracle is precisely what Agile-SOFL workflows need.  Agile development
cycles are short (two-week sprints); full formal verification on every candidate is
impractical.  The Z3 checker provides meaningful formal grounding — the generated test
cases are derived directly from the FSF specification structure, not from manually written
unit tests — at a cost compatible with sprint cadences.  The +3.9\,pp conformance
contribution of the checker (A1 ablation) demonstrates that this grounding materially
improves output quality, not merely adds bureaucratic overhead.

\subsection{Pattern Guard Selectivity and Advisory Use}

The A2 ablation result is the most nuanced finding in the evaluation.  Removing the pattern
guard yields a marginal \emph{improvement} of $+0.5$\,pp in formal conformance, suggesting
that the guard does not improve pass rates at the current benchmark scale.  At first glance
this appears to undermine the case for including the pattern guard in the pipeline.

A more careful interpretation, however, is that the pattern guard provides value as a
\emph{classification and diagnostic tool} rather than as a \emph{filtering gate} that
improves overall pass rates.  When the guard fires, it identifies the specific fault pattern
involved (e.g.\ RF07, unconditional success; or RF03, missing precedence guard), providing
actionable information to the developer.  This diagnostic value is not captured by the
conformance metric, which measures only whether test cases pass, not \emph{why} they fail.

We therefore recommend a modified deployment strategy: the pattern guard should emit
\emph{advisory warnings} rather than hard rejection signals for most fault patterns,
following static-analysis guidance that severity should modulate gating policy~\citep{hovemeyer2004findbugs,ayewah2008findbugs}.
The exception is RF07 (unconditional success), a pattern that trivially satisfies test cases
while violating all meaningful specification constraints; for RF07, hard rejection remains
appropriate.  This approach preserves the diagnostic benefit of the pattern guard while
avoiding the marginal conformance penalty incurred by its conjunctive gating in the current
configuration.

% ---------------------------------------------------------------------------
\section{Limitations}
\label{sec:disc:limitations}
% ---------------------------------------------------------------------------

Several limitations bound the scope of the claims made in this thesis.

\paragraph{Benchmark representativeness.}
The hard benchmark consists of 120 synthetic Agile-SOFL tasks generated from a fixed
template family.  While the tasks were specifically designed to probe precedence-sensitive
scenario orderings — the class of faults HSP-Agile targets — they may not capture the full
diversity of FSF specifications encountered in industrial practice.  Industrial
specifications may exhibit more complex state dependencies, external service interactions,
or composite scenario structures not represented in the benchmark.  Replication on
industrial specifications from real Agile-SOFL projects~\citep{li2022agilesofl,pohl2010requirements} is necessary before broad
generalisability can be claimed.

\paragraph{Single LLM family.}
All experiments used Qwen2.5-27B-Instruct via an OpenAI-compatible LLM gateway.  LLMs differ substantially in their tendency
to produce scenario-ordering violations, their responsiveness to repair prompts, and their
latency characteristics.  The conformance gains reported here may be larger or smaller with
GPT-4, Claude~3.5 Sonnet, or Gemini~1.5 Pro.  The qualitative structure of the pipeline
(checker-guided repair) is model-agnostic, but the quantitative results are specific to the
evaluated model family.

\paragraph{Z3 test-case coverage.}
The formal checker currently generates one test case per scenario branch.  Generating
multiple cases per branch — for example, using boundary-value analysis or random
perturbation around the Z3 witness — would increase fault detection sensitivity but also
increase latency.  The current single-case-per-scenario regime likely underestimates the
conformance advantage of mode~M relative to what a denser test suite would reveal.

\paragraph{Prevention evaluation completeness.}
The probabilistic detection rate (PDR) and false accept rate (FAR) analysis reported in
Chapter~\ref{ch:experiments} is based on 40 tasks per mode.  Preliminary figures
(B1: 7.5\%, B2: 15.0\%, M: 10.0\% PDR) are consistent with conformance trends but require
a full mutation-injection experiment across all modes and larger mutant populations before
definitive prevention claims can be made.  The current figures should be treated as
indicative rather than conclusive.

\paragraph{Repair loop convergence.}
The maximum repair budget $K$ is fixed across all tasks.  Harder tasks may benefit from
larger $K$, while simpler tasks waste budget on unnecessary iterations.  An adaptive
budget-allocation strategy that dynamically scales $K$ based on observed conformance
improvement per iteration could reduce latency while maintaining quality on the hardest
tasks.

% ---------------------------------------------------------------------------
\section{Comparison to Related Approaches}
\label{sec:disc:related}
% ---------------------------------------------------------------------------

\paragraph{DafnyBench and full deductive verification.}
DafnyBench~\citep{loughman2024dafnybench} evaluates LLM synthesis of Dafny programs with machine-checked
correctness proofs.  Full deductive verification provides stronger guarantees than
HSP-Agile's bounded Z3 checking: a Dafny proof is valid for all inputs, whereas
HSP-Agile's conformance certificate covers only the generated test suite.  However,
Dafny's verification requires formal annotations (invariants, pre/post-conditions, loop
measures) that are absent from FSF specifications.  HSP-Agile is applicable directly to
Agile-SOFL FSF inputs without requiring the developer to write verification annotations,
making it usable by practitioners who are not trained in deductive verification.  The two
approaches occupy complementary positions in the verification-cost spectrum.

\paragraph{CEGIS-inspired approaches for Dafny (Lynn et al., TOSEM 2024).}
Lynn et al.~\citep{lynn2024verifierloop} apply counterexample-guided inductive synthesis (CEGIS)
to assist LLMs in producing correct Dafny programs.  The conceptual similarity to
HSP-Agile's repair loop is strong: both approaches use witnessed failures to construct
targeted repair prompts.  The key difference is that Lynn et al.\ operate in the Dafny
verification context, where counterexamples are generated by the Dafny verifier from
logical proof obligations, whereas HSP-Agile generates counterexamples from Z3-driven test
cases derived from FSF scenario semantics.  HSP-Agile's checker is lighter-weight and does
not require a verification-capable specification language, but it provides weaker
guarantees as a consequence.

\paragraph{Standard LLM baselines (B1 and B2).}
Relative to the two LLM baselines evaluated in this thesis, HSP-Agile adds
$+6.2$\,pp and $+2.4$\,pp formal conformance respectively, at a latency cost of
approximately $4.6\times$ over B2~\citep{dakhel2023copilot,chen2021codex}.  These figures place HSP-Agile in a ``high quality,
moderate cost'' regime: it is not suitable as a real-time code-completion assistant but is
appropriate for quality gates in workflows where downstream formal verification costs
dominate.
